% from ega4-12-fr.tex
\oldpage[IV-3]{178}
propriété $(S_k)$ dans un voisinage $U$ de $y$ dans $Y$ ((6.11.2) et (7.8.3,
(iv)))~; on conclut alors de la même manière en remplaçant cette fois $Y$ par
$U$ et $X$ par $V\times_YU$.
\end{proof}

\begin{theorem}[12.1.6]
\label{IV.12.1.6-fr}
Soient $f:X\to Y$ un morphisme plat et localement de présentation finie, $k$
un entier. Les parties suivantes de $X$ sont ouvertes~:
\begin{enumerate}[label=(\roman*)]
  \item L'ensemble des points $x\in X$ tels que $X_{f(x)}$ vérifie la propriété
  $(S_k)$ au point $x$.
  \item L'ensemble des $x\in X$ tels que $X_{f(x)}$ vérifie la propriété
  $(R_k)$ géométrique au point $x$, soit équidimensionnel au point $x$ et tels
  en outre que $x$ n'appartienne à aucun cycle premier immergé associé à
  $X_{f(x)}$.
  \item L'ensemble des $x\in X$ tels que $X_{f(x)}$ soit géométriquement
  régulier (i.e. lisse) au point $x$.
  \item L'ensemble des $x\in X$ tels que $X_{f(x)}$ soit géométriquement
  normal au point $x$.
\end{enumerate}
\end{theorem}

\begin{proof}
Les étapes A) et B) de (12.0.2) se font ici comme dans (12.1.1)~; pour l'étape
A), on utilise (6.7.8), ainsi que (4.2.7) pour (ii)~; pour l'étape B), on
utilise (9.9.2) et (9.9.4). Reste à examiner l'étape C) dans chaque cas.
\begin{enumerate}[label=(\roman*)]
  \item[(i)] Comme dans (12.1.1.2), on se ramène (en utilisant (6.7.8)) au cas
  où $Y$ est le spectre d'un anneau de valuation discrète $A$,
  $X=\operatorname{Spec}(B)$, où $B$ est une $A$-algèbre de type fini, $x$,
  $x'$ deux points de $X$ tels que $f(x)=y$ soit le point fermé et
  $y'=f(x')$ le point générique de $Y$, $x'$ étant en outre une générisation
  de $x$. Comme il s'agit de prouver que $X$ vérifie $(S_k)$ au point $x'$,
  on peut remplacer $X$ par $\operatorname{Spec}(\mathcal O_{X,x})$,
  c'est-à-dire supposer l'anneau $B$ local, l'homomorphisme $A\to B$ étant
  local et $B$ étant une $A$-algèbre essentiellement de type fini (1.3.8).
  Alors, par hypothèse, si $t$ est une uniformisante de $A$, $t$ est un
  élément $B$-régulier par platitude, et $B/tB$ vérifie $(S_k)$. Mais comme
  $A$ est un anneau universellement caténaire (5.6.4), $B$ est caténaire,
  donc, par (5.12.4), il en résulte que $B$ vérifie $(S_k)$, ce qui achève la
  démonstration dans le cas (i).
  \item[(iii)] Ici l'étape C) est inutile~; comme $f$ est plat, on sait en
  effet (6.8.7) que lorsque $Y$ est localement noethérien et $f$ localement de
  type fini, l'ensemble des $x\in X$ tels que $X_{f(x)}$ soit géométriquement
  régulier au point $x$ est ouvert dans $X$.
  \item[(ii)] En raisonnant comme dans (12.1.1.3), pour prouver que la propriété
  considérée est stable par générisation, on peut d'abord, en considérant une
  extension finie quelconque $K'$ de $k(y')$, et tenant compte de la définition
  (6.7.6) de la propriété $(R_k)$ géométrique, ainsi que de l'invariance par
  changement du corps de base des deux autres propriétés figurant dans (ii),
  remplacer dans (ii) la propriété $(R_k)$ géométrique par la propriété
  $(R_k)$. Procédant alors comme dans (i), on se ramène au cas où $Y$ est
  spectre d'un anneau de valuation discrète $A$, et puisque $A$ est caténaire,
  il suffit d'appliquer (5.12.5) pour conclure.
  \item[(iv)] L'ensemble des points $x\in X$ tels que $X_{f(x)}$ soit
  géométriquement normal au point $x$ est contenu dans celui des $x\in X$ tels
  que $X_{f(x)}$ soit géométriquement ponctuellement intègre et vérifie $(S_2)$
  au point $x$, et ce dernier est ouvert dans $X$ en vertu de (12.1.1,
  (viii)). On peut donc déjà supposer que, pour tout $x\in X$, $X_{f(x)}$ est
  géométriquement ponctuellement intègre et vérifie $(S_2)$ au point $x$, et
  \emph{a fortiori} il est équidimensionnel. D'autre part, en vertu du critère
  de Serre (5.8.6), dire que $X_{f(x)}$ est géométriquement normal au point $x$
  signifie que $X_{f(x)}$ vérifie $(S_2)$ et la propriété $(R_1)$ géométrique
  en $x$~;
  